\documentclass[10pt]{article}

\usepackage[utf8]{inputenc}
\usepackage[spanish]{babel}
\usepackage{setspace}

\usepackage{graphicx}
\usepackage[shortlabels]{enumitem}
\usepackage{url}
\usepackage[margin=1in]{geometry}

\newcommand{\subsubsubsection}[1]{\paragraph{#1}\mbox{}\\}
\setcounter{secnumdepth}{4}
\setcounter{tocdepth}{4}


\title{
    \includegraphics[width=0.3 \linewidth]{Logos/UIS.pdf} \\
    Objetividad: Un argumento para obligar}
\author{Daniel David Delgado Cervantes \\  
        Laura Alexandra Hernández Pérez\\ 
        Grupo \# 4
}

\onehalfspacing

\begin{document}
    
\maketitle

    \section{Introducción}

    Una de las costumbres occidentales, en cuanto se refiere a la necesidad de hacer que alguien concuerde con alguna idea o concepto, sin el uso de la fuerza, es la búsqueda de un argumento sin ningún margen de discusión. Este tipo de argumentos, comúnmente conocidos como argumentos objetivos, son normalmente aquellos que se basan directamente con la realidad y están basados en hechos. A pesar de esto, Humberto Maturana cuestiona esta definición de objetividad debido a que, en términos simples, este concepto niega, o por lo menos no reconoce, las propiedades biológicas que los observadores, en especial los humanos, tenemos al ser sistemas vivientes.

    Es a partir de esto de este cuestionamiento sobre el concepto de la objetividad trabajándolo desde el reconocimiento de tanto los sentidos como fenómenos biológicos mutables que Maturana desarrolla los conceptos de objetividad al igual que lo que entendemos como la percepción desde el punto de vista para una entidad biológica, es decir, la percepción de la realidad desde un punto de vista filosófico y biológico. 

    \section{La ontología del explicar}

    El principal objetivo de este capítulo está orientado a la búsqueda del entendimiento de los conceptos relacionados con los sucesos de nuestra diario vivir al igual que las explicaciones que usamos para poder entender el mundo que nos rodea al igual que nos entendemos a nosotros mismos.

        \subsection{Praxis Del Vivir}

        El concepto de la praxis del vivir es uno de los más prominentes dentro del discurso de Maturana. Esto se debe a que los que se considera la praxis del vivir, que puede tomarse de manera casi literal, es el concepto de nuestro diario vivir o lo que nosotros entendemos como las vivencias de cada día.
        
        De la misma manera, este concepto puede ser tomado como el como nosotros, en rol de observadores, interpretamos la vida a través de los sentidos y formulamos explicaciones con el fin de entender el mundo. En este sentido, como ya hemos establecido, las explicaciones que realizamos para entender el mundo, son las que determinan el como definimos nuestra propia praxis del vivir. 

        \subsection{Explicaciones}

        Como hemos establecido, las explicaciones son una parte importante del como interpretamos el mundo y del como entendemos la praxis del vivir. Esto se debe a la función que cumplen de manera tradicional en tanto estas son usadas para responder preguntas. Cuestionamientos relacionados con sucesos normalmente búscan una explicación con el fin de resolver las dudas que el observador tenga respecto a las eventualidades que surjan dentro su praxis del vivir. En este sentido, es necesario entender la función que cumplen las explicaciones en nuestro rol de observadores a medida con observamos el mundo a través de los sentidos. 

        Lo primero a resaltar de las explicaciones es el origen de estas. En términos generales, entendemos que las explicaciones son el mecanismo a través del cual nosotros, en rol de observadores, entendemos los eventos que toman lugar dentro de nuestra praxis del vivir. Es decir, las explicaciones surgen del observador a partir de las eventualidades que este experimenta. Es de esto donde se entiende el porque las explicaciones son consideradas de segundo orden. Ya que las explicaciones se crean desde los sucesos, estas pueden ser consideradas \textit{innecesarias} puesto que no afectan en medida alguna el suceso.  

        De igual manera, estas explicaciones están a merced del criterio de aceptación del observador en tanto este es el responsable de determinar qué tipo de explicaciones son aceptadas dentro de su praxis del vivir. Es decir, el criterio de aceptación es el filtro el cual determina qué explicaciones, sean formuladas por el observador o algún interlocutor, desembocan en una reformulación de la praxis del vivir, o en un cambio en la manera en la que el observador comprende su realidad.

        \subsection{Caminos Explicativos}

        Uno de las primeros planteamientos necesarios para comprender los diferentes caminos explicativos está dentro del como, los que toman el rol de observador, perciben el mundo a través de sus habilidades cognitivas, al igual que estos reflexionan sobre dichas habilidades. 

        El primero de los dos caminos explicativos, llamado el camino de la objetividad sin paréntesis, u objetividad trascendental; hace referencia a una percepción del mundo en la cual el observador y la realidad son independientes uno del otro. Con esto Maturana se refiere a una visión del mundo en la cual la interpretación y las vivencias no dependen del observador, puesto que este asume que la realidad que percibe es veráz y no posee márgen de discusión en cuanto lo que sucede, es percibido tal y como es.

        En cuanto al segundo camino explicativo, es aquí donde el observador comprende que él, como sistema viviente, posee ciertas habilidades cognitivas que a su vez son fenómenos biológicos. De igual manera, el observador comprende que estos fenómenos biológicos pueden verse afectados como cualquier otro sistema. En este sentido, el observador está dentro de la objetividad en paréntesis, o objetividad constructivista; en la cual este entiende las limitaciones de sus sentidos en tanto estos no pueden diferenciar lo que es una percepción o lo que es un engaño. Es decir, el observador, en su compresión de los sentidos como fenómenos biológicos, entiende que lo que percibe, depende de él. O dicho de otra manera, el observador entiende que su realidad es construida por si mismo con el uso de sus sentidos. 

        La principal consecuencia de este tipo de objetividad, está relacionado con el entendimiento de las realidades de otros observadores. Es decir, para los observadores que viven dentro de este camino de la objetividad entienden que para una distinción realizada sobre un evento o suceso, aunque puede ser completamente distinta a la realizada por él, es igualmente válida en todo sentido. Esto se debe a que, para todos los observadores, sus distinciones propias, tienen un coherencia dentro de sus praxis del vivir debido a que cumplen con el criterio de aceptación.

        \subsection{Dominios Explicativos}

        Los dominios explicativos es uno de los conceptos base en cuanto a la compresión del como entendemos la realidad. Como su nombre lo indica, un dominio explicativo se refiere al conjunto de explicaciones, que tanto realizamos como recibimos en nuestro rol de observadores, las cuales consideramos como válidas. 
        
        A partir de esta descripción, es posible inferir que el dominio explicativo está compuesto de lo que conocemos como la praxis del vivir al igual que el criterios de aceptación que tengamos para la validación de las explicaciones que vamos formulando dentro de nuestro diario vivir.  

        Ahora, el concepto más importante dentro de los dominios explicativos recae en la posibilidad de llegar a consensos de tanto criterios de aceptación como de maneras en las cuales formulamos explicaciones. Es decir, a partir de los dominios explicativos, podemos llegar a la generación de consensos los cuales darían como resultado explicaciones universales para cualquier fenómeno que un observador explique. O en otras palabras, la posibilidad de tener las mismas explicaciones con observadores completamente diferentes.

        Este aspecto es especialmente útil en cuanto tenemos la capacidad como observadores de trabajar bajo una misma estándar de realizar las cosas. Uno de los casos en los que esto es especialmente útil, es en la generación de explicaciones científicas. En este sentido, la capacidad de mantener constantes las maneras en las que se realizan las explicaciones, permite la universalidad en el conocimiento o por lo menos en el como entendemos el funcionamiento de las cosas.

    \section{Realidad: Una Proposición Explicativa}

    El capítulo sobre la realidad tiene como objetivo principal el desarrollo de las consecuencias que tienen los caminos explicativos dentro de algunos de los aspectos que hacen parte de nuestro diario vivir. Estos temas van desde el lenguaje hasta lo que consideramos como el ser consientes.

        \subsection{Lo Real}

        
        
        \subsection{Racionalidad}

        
        
        \subsection{Lenguaje}



        \subsection{Emotividad}


        \subsection{Conversaciones}

        Las conversaciones, en términos generales, comprende el como fluye el lenguaje entre observadores con el fin de coordinar acciones y emociones con el fin de realizar o cumplir ciertos objetivos. Para explicar algunos de los tipos de conversaciones que toman lugar dentro de nuestra praxis del vivir, Maturana da los siguientes ejemplos:

        \begin{enumerate}[a.]
        \item Conversaciones de coordinaciones de acciones presentes y futuras.
        
        Consisten en conversaciones que se componen únicamente de coordinaciones de acciones.

        Ejemplo: ``Si tu pones la mesa, yo preparo la cena. / Lo haré con placer.''

        
        \item Conversaciones de quejas y disculpas por acuerdos no tomados.
            
        Básicamente se refiere a faltas o fallas a conversaciones de coordinaciones de acciones.
                
        Ejemplo: ``Estoy listo ahora. ¿Estás listo? / Lo siento, no puedo hacerlo ahora. / Pero me prometiste / Sí, pero mi madre me está llamando, ¿puedes esperarme?''
        
        \item Conversaciones de deseos y expectativas.
            
        Se refiere a conversaciones donde se presenta la descripción de un futuro ignorando las acciones que constituyen a la persona en el presente.

        Ejemplo: ``Después de la elección presidencial seré capaz de impulsar mi programa de reforestación. / Ese será el caso si tu candidato gana. Yo pienso que sin embargo que él no lo hará. / Yo estoy seguro que él ganará; él tiene el respaldo de los trabajadores.''
        
        \item Conversaciones de mandos y obediencias.

        Se refiere a conversaciones donde una figura de autoridad doblega a alguien para que realice una acción la cual este no se encuentra de acuerdo. El resultado de esto es el ceñir con más fuerza la posición de autoridad del que se encuentra al mando.
                
        Ejemplo: ``Juan, ven a resolver este problema en el pizarrón. / Pero no he terminado aún el ejercicio en mi cuaderno. / No importa, te estoy pidiendo que vengas al pizarrón.''
        
        \item Conversaciones de caracterizaciones, atribuciones y evaluaciones.
            
        Estas conversaciones se refieren, de manera general, a aquellas en las cuales se realizan descripciones y opiniones sobre las percepciones del espectador. 
                
        Ejemplo: ``¡Aquí estás! Yo te creía una persona que siempre llegaba a la hora. / ¿Qué?, ¿Quieres decir que soy impuntual? Esta es la primera vez que me atraso.''
        
        \item Conversaciones de quejas por expectativas incumplidas
            
        Estas se relacionan con los casos donde el oyente es culpado del rompimiento de una promesa realizada con anterioridad sobre el cumplimiento de algo prometido.
                
        Ejemplo: ``Tu llegas tarde de nuevo y la comida está recocida. / ¡Pero tu sabes que en esta época del año no puedo llegar más temprano!''

        \end{enumerate}

        \subsection{El sistema nervioso}

        Como primera medida, Maturana resalta el aspecto principal del sistema nervioso de manera general. Él lo describe como una red cerrada de componentes que interactúan entre sí dentro de un sistema de mayor tamaño. En nuestro caso, este sistema nervioso está compuesto de neuronas, células sensoriales y células efectoras. 

        De la misma manera, resalta el determinismo estructural (Concepto introducido por Maturana que establece que las perturbaciones que puede sufrir un sistema al interactuar con el ambiente están limitadas por la dinámica de interacciones que le permitan su estructura únicamente. El agente ambiental que efectúa la perturbación no tiene nade que ver en ese aspecto.) de nuestro sistema nervioso no nos permite distinguir entre la percepción y la ilusión. 

        Como segundo punto se establece que, el resultado de las interacciones del sistema nervioso con en ambiente. Esto se debe a que, el sistema nervioso, como parte de un ente vivo, a medida que pasa el tiempo y dependiendo de las interacciones que este perciba con el ambiente, puede modificar su estructura con los componentes.

        El tercer punto se refiere al como, en la medida en la que hay cambios a la estructura debido a sus interacciones con el ambiente, da como resultado un cambio de estado en el organismo el cual afecta sus interacciones. De esto se concluye que el sistema nervioso de un organismo es el presente de un organismo.

        Finalmente, Maturana determina que, dadas las condiciones en las que, un organismo participa en un dominio del lenguajeo, este puede puede operar dentro de un flujo de dinámicas internas, llamadas diálogos internos. Esto se generó a partir del sistema nervioso el cual le permitió participar del lenguaje. 

        \subsection{Autoconciencia}

        Autoconciencia, en términos generales, es lo que nosotros conocemos con la realización de una distinción de un sistema viviente sobre si mismo. Es decir, la autoconciencia es tener el conocimiento sobre la relación entre uno mismo y el entorno que lo rodea. 

        En este sentido, es necesario denotar que nosotros, como seres humanos, realizamos este ejercicio de ser autoconscientes a través del mecanismo que conocemos como el lenguaje. Esto se debe a que, como es de esperarse, el fenómeno de la autoconciencia requiere de las reflexiones y las explicaciones que podemos realizar a través del lenguaje con el fin de realizar preguntas sobre nosotros mismos. Como es posible apreciar, esto tiene una fuerte relación con el sistema nervioso en tanto este permitió los diálogos internos.
        
        De igual manera, este fenómenos empieza en el momento en el que realizamos estas preguntas y buscamos las explicaciones sobre las distinciones de nuestro cuerpo respecto al de otros observadores. En las palabras de Maturana, "Todo el dominio de la autoconciencia surge como un dominio de recursión en el darse cuenta de uno mismo.".

        \subsection{Epigénesis}

        Los conceptos de la epigénesis están relacionados con lo que conocemos con la estructura del cuerpo y la capacidad del mismo en cuanto a estos determinan lo que este puede y no puede hacer.  Esto se refiere que para cualquier sistema viviente, este está limitado dentro de su estructura inicial. Esto se da a partir de lo que conocemos como el determinismo genético, el cual está definido como la especificación de los ácidos nucléicos de un futuro en el desarrollo de un organismo. En otras palabras, nuestro ADN define lo que podemos y no podemos hacer dentro de nuestra historia.

        De igual manera, Maturana se refiere a la distinción que debe realizarse al aprendizaje. Esto se debe a la mala concepción del aprendizaje como una ``adaptación'' al medio ambiente al cual este pertenece. Este problema de la adaptación es tratado desde los dos caminos explicativos. En el camino explicativo de la objetividad sin paréntesis, según Maturana, se dice que el aprendizaje es un procesos dirigido por la adaptación. Básicamente, el proceso de aprendizaje se ve como un comentario sobre dos momentos de la epigénesis en el cual no se reconoce el proceso histórico y se asume que existe un mecanismo activo de acomodación. Por otro lado, en la objetividad en paréntesis, el aprendizaje es un proceso epigénico puesto que no implica la construcción de una representación del medio ambiente.

    \section{Ontología del Conocer}

        \subsection{Observador - Observación}

        \subsection{Conocer}

        \subsection{Interacciones Mente y Cuerpo}

    \section{Lo social y Lo Ético}

        \subsection{Lo social}

        Como punto de partida, se nos introducen los conceptos del amor y Para Maturana el concepto del amor es la emoción en la cual las interacciones sociales o relaciones sociales existen. En este sentido, Maturana no se refiere al amor como bondad sino como el fenómeno biológico en el cual los sistemas vivientes pueden coordinar sus acciones como consecuencia de la aceptación mutua. En este sentido, se sostiene que los sistemas sociales son constituidos por sistemas vivientes los cuales, bajo la emoción del amor, desarrollan coordinaciones de acciones de aceptación.

        En este sentido, los sistemas sociales, trabajando bajo la emoción del amor, tienen algunas consecuencias en cuanto al como se definen estos:
        
        \begin{enumerate}[a)]
            \item La primera consecuencia de esto está en que los sistemas sociales, de una u otra manera, requieren la existencia de sistemas vivos para poder existir. En este sentido, en el caso de que se eliminen o nieguen los sistemas vivos, se destruirían o eliminarían los sistemas sociales.   
            
            \item La segunda se relaciona con el concepto de identidad de clase en el sentido que la identidad define el sistema social. Es decir, un sistema social compuesto de sistemas vivos cuya identidad sea la de un ser humano, dará como resultado un sistema social de seres humanos. De la misma manera, cualquier cosa que niegue la identidad de los componentes, destruye el sistema social.
            
            \item La tercera se refiere al medio por el cual se da este sistema social, es decir, a través de donde se da este mecanismo social. En el caso de las personas, nuestros sistemas sociales, como lo afirma Maturana, surgen únicamente del lenguajeo. 
            
            \item La cuarta está relacionada con los sistemas vivientes y sus interacciones con otros sistemas vivientes. Maturana define que, para un observador, un sistema viviente hace parte de un sistema social si este interactúa con otros sistemas vivientes pertenecientes a dicho sistema social.
            
        \end{enumerate}

        De igual manera, se nos presenta las condiciones necesarias para la conservación de un sistema social. La principal necesidad de un sistema social se relaciona con las interacciones de los mismos miembros de sistema social y su capacidad de trabajar dentro de la conducta del sistema social. En el caso donde los miembros dentro del sistema social actúen por fuera de la búsqueda de la aceptación mutua entre sus miembros, el sistema social se desintegrará.

        Finalmente, se presentan los conceptos de los límites de los sistemas sociales. Lo primero a resaltar está dado por el como, los límites están dados únicamente por los sentimientos de los integrantes del sistema. En este sentido, esto puede explicarse de dos maneras. En la objetividad sin paréntesis, se dice que la razón de estos límites está basado en los argumentos racionales usados para rechazar o aceptar a los nuevos integrantes de un sistema social. Es decir, se ignoran las emociones. Por el contrario, en la objetividad en paréntesis, se tienen en cuenta que las emociones hacen parte fundamental de la aceptación o negación de los miembros a ingresar dentro de un sistema social. O dicho de otra manera, las barreras de un sistema social sólo podrán ser atravesadas a partir de la seducción emocional.

        \subsection{Multiplicidad de los Dominios de Existencia}

        En este capítulo, se hace especial relevancia hacia los conceptos de los sistemas no sociales. A diferencia de los sistemas sociales, este tipo de sistemas aunque trabajan bajo las coordinaciones de acciones de los integrantes del sistema social, la razón de las coordinaciones de acciones no se dan desde el amor.

        Los dos sistemas no sociales más relevantes son el sistema laboral, donde las coordinaciones de acciones se dan para el cumplir con un objetivo; y el sistema jerárquico, en el cual las coordinaciones de acciones se dan a partir de la necesidad de cumplir una orden dada. Como es de denotar, estos sistemas funcionan con base en el respeto donde la emoción de aceptación está más relaciona con la obediencia hacia las órdenes dadas. De estos, también es posible resaltar algunas características:
        
        \begin{enumerate}[a)]
            
            \item No se intersectan con otras comunidades de manera directa pero nuestras interacciones con estas sí cambia el como nos comportamos respecto a todo.
            
            \item Los sistemas vivientes que interactúan con las comunidades, en términos simples, son el resultado de la historia de las interacciones dentro de las comunidades.
            
            \item Los cambios dentro de las comunidades sólo ocurre como un cambio de la corporalidad de los miembros de la comunidad.
        \end{enumerate}

        \subsection{Lo ético}
        
        Lo primero para entender el porque la ética es meramente emocional recae en las emociones, está dado en entender que las emociones son un fenómeno biológico propio de la corporalidad. De la misma manera, se debe entender que la cultura no afecta nuestras emociones pero sí el emocionar.
        
        De esto, se presenta la relación con las consideraciones éticas y el porque estas se realizan. Maturana resalta que este tipo de planteamientos sólo se realizan cuando existe un quiebre en el respeto humano y se buscan que hayan consecuencias hacia un grupo en particular. Es decir, la ética viene de las emociones y no de la racionalidad. Con esto se refiere que depende del como se vea un problema y que tal le parezca al observador. Si le parece bien, no se hará la distinción; de lo contrario, se hará.

        Se resaltan los límites de interés del bienestar de los seres humanos los cuales están dados por la ubicación de otro ser humano dentro de los distintos sistemas sociales. Esto se refiere a los casos donde realmente no nos interesa el bienestar de otro ser humano debido a su condición de estar dentro de otro sistema social que puede que el observador no considere válido dentro de su praxis del vivir. Es de denotar que cuando un observador externo acepta una posición o conducta ética nuestra, no es que el observador acepte como racional y por ende objetivo argumento, sino que entiende y acepta el dominio existencial trabajado. Finalmente se determina que, a medida que vamos moviéndonos dentro los diferentes dominios sociales, sean sociales o no sociales, nuestra preocupación va cambiando dependiendo de nuestras emociones.
        
        Estos límites pueden ser explicados de dos maneras. En el camino de la objetividad sin paréntesis entra en conflictos debido a que se está intentando encontrar una validez simultanea de las emociones y de la razón, cosa que no va a pasar. Por el contrario, en la objetividad en paréntesis, debido a que como sistemas vivientes entendemos que podemos vivir en diferentes dominios de existencia, podemos trabajar bajo nuestros fundamentos éticos al igual que se entiende que los intereses éticos van hasta nuestras fronteras operacionales. 

\end{document}